\section{Continuous Proof: From Changed Paths to Merge Witness}

Agent-native testing should be organized by proof question, not author convenience. Domain invariants need unit and property tests. Application workflows need authorization, idempotency, and transaction tests. Adapters need integration and contract tests. Durable data stores need migration, constraint, and schema drift checks. The browser needs component tests, rendered UX proof, and Playwright coverage for critical paths. Public benchmarks and open-source agents show that setup reliability, context retrieval, and interface design are central bottlenecks for autonomous repair \cite{sweAgentAci2024,sweBench2024,setupBench2025,contextBench2026,openhands2024,aiderGithub2026}.

\begin{table}[t]
\caption{Minimum proof lanes by ownership cell.}
\label{tab:proof-lanes}
\centering
\small
\begin{tabularx}{\columnwidth}{L{0.28\columnwidth} Y}
\toprule
\JKTableHeader
\textbf{Layer} & \textbf{Minimum proof} \\
\midrule
Domain & Unit tests plus property tests for invariants and state machines. \\
Application & Integration tests for authz, idempotency, transactions, workflows. \\
Adapters & DB/external integration tests and failure injection. \\
Contracts & Generation check, drift check, compatibility check. \\
Durable data & Migration apply, constraint tests, RLS or tenant-policy tests where used. \\
Web & Typecheck, component tests, rendered UX QA, Playwright critical paths. \\
AI/data service & Contract tests, eval suites, no direct product-truth ownership. \\
Ops/security & Secret scan, dependency scan, provenance, audit gate. \\
\bottomrule
\end{tabularx}
\end{table}

The explosion of tests is useful only if routed. A repo with thousands of tests and no \texttt{test-map.json} still asks agents to guess. The standard therefore requires a fast lane for changed files, a contract lane for generated interfaces, a security lane for dependency and secret risk, a DB lane for migrations, a web lane for component and rendered UX proof, an e2e lane for critical browser behavior, and a release lane for full confidence.

Refactoring is safe only when changed behavior has routed proof. A rename, extraction, deletion, or boundary move should be rewarded when it reduces debt, but it still needs receipts proving that the owner route, generated boundaries, and behavior claims remain valid. Agents otherwise tend toward additive patches: new wrappers, new tests beside old ones, new adapters beside stale paths. A Jankurai repository should make deletion, narrow repair, and proof-backed refactoring cheaper than piling on another plausible layer.

\begin{lstlisting}[style=codebox]
[[lane]]
name = "web"
command = "just ux-qa"
cwd = "."
timeout_seconds = 120
network = "loopback_only"
write_paths = ["target/jankurai/ux/"]
artifacts = [
  "target/jankurai/ux/evidence.json",
  "target/jankurai/ux/screenshots/"
]
\end{lstlisting}

CI modes should be explicit. \emph{Advisory} mode emits reports but does not fail merge. \emph{Guarded} mode blocks critical caps. \emph{Standard} mode enforces high and critical findings plus required receipts. \emph{Ratchet} mode prevents score or finding regression without a named waiver. \emph{Release} mode gates shipped artifacts on audit, contracts, security, DB, e2e, versions, and release evidence.

The continuous proof loop is concrete and local. A changed path enters \texttt{agent/owner-map.json} for ownership and \texttt{agent/test-map.json} for lane routing. \texttt{jankurai lane} or \texttt{jankurai proof} turns that route into a proof plan. \texttt{jankurai prove} executes allowed commands, writes command receipts and \texttt{target/jankurai/evidence-index.json}, and records artifact digests. \texttt{jankurai proof-verify} checks the plan and evidence against the current repository state. \texttt{jankurai audit} folds ownership, generated zones, proof receipts, security evidence, UX evidence, and score history into JSON/Markdown/SARIF/JUnit reports. \texttt{jankurai witness} compares the candidate against an accepted baseline and emits the PR proof matrix. Findings become a repair queue through \texttt{agent\_fix\_queue}, \texttt{issues export}, or \texttt{repair-plan}. This is the real-time part of Jankurai: not an omniscient daemon, but a repeatable local/PR/CI loop that makes merge evidence current.

\begin{table}[t]
\caption{Proof freshness algorithm.}
\label{tab:proof-freshness}
\centering
\scriptsize
\begin{tabularx}{\columnwidth}{L{0.26\columnwidth} Y}
\toprule
\JKTableHeader
\textbf{Check} & \textbf{Invalid when} \\
\midrule
Commit binding & Receipt \texttt{git\_head} differs from witness head. \\
Changed path set & Receipt omits a changed path in its declared lane or covers paths no longer in the witness. \\
Routing maps & Owner map, test map, or generated-zone map changed after the receipt was produced. \\
Lane authority & Command is absent from the current proof lane map. \\
Tool identity & Tool version, command digest, standard version, auditor version, or schema version is missing or mismatched. \\
Artifact digest & Named log, report, screenshot, or trace is missing or hash-mismatched. \\
Dirty state & Worktree is dirty when the lane, witness, or release policy requires a clean tree. \\
Privacy status & Required redaction status is absent for restricted evidence. \\
\bottomrule
\end{tabularx}
\end{table}

Flaky proof is not proof. A lane may retry only under a declared retry policy that records attempts, timing, environment, and failure class. A pass-after-retry can satisfy advisory evidence, but release and ratchet gates should require either deterministic rerun stability or a waiver with owner approval, residual risk, and a repair plan for the flaky lane.

\begin{lstlisting}[style=codebox]
jankurai witness . \
  --changed-from origin/main \
  --baseline agent/repo-score.json \
  --out target/jankurai/merge-witness.json \
  --md target/jankurai/merge-witness.md
\end{lstlisting}

The report itself is proof material. JSON and Markdown are the agent repair surface; SARIF, JUnit, GitHub step summaries, JSONL repair queues, and issue exports are integration surfaces. A failed score with no findings is invalid output: every below-floor dimension must produce at least one routed finding with owner, lane, rerun command, evidence kind, fingerprint, and queue item.
